\subsection{Линейные уравнения}

\begin{warn}
Название главы взято из тетради дословно: «Уравнения в частных производных
первого порядка». У первой главы названия нет ни на доске, ни в тетради —
изложение там начинается сразу с \S1, так что заголовок «Волновое уравнение»
придуман для этого конспекта.
\end{warn}

Однородное уравнение в частных производных первого порядка:

\stepcounter{equation}
\begin{equation}\tag{1}\label{eq:fo-homog}
  \sum_{i=1}^{n} a_i(x)\, u_{x_i} \equiv \bar a(x) \cdot \nabla u = 0 ,
\end{equation}
где коэффициенты $a_i$ гладкие и $\bar a \ne 0$.

\inputfig{figures/fig-char-field}

\subsubsection*{Пример: уравнение переноса}

\begin{example}
Среда (жидкость, газ) движется параллельно оси $x$ с постоянной скоростью $a$.
Пусть $u(x,t)$ — концентрация частиц (числовая плотность, то есть число частиц
на единицу длины).

\medskip
\textbf{Закон сохранения.} Скорость изменения числа частиц на отрезке
$[x_1,x_2]$ равна разности втекающего и вытекающего потоков:
\[
  \dd{}{t} \int\limits_{x_1}^{x_2} u(x,t)\, dx = j(x_1,t) - j(x_2,t),
\]
где $j$ — поток частиц. Для среды, которая просто сносит частицы со скоростью
$a$,
\[
  j(x,t) = a\, u(x,t) .
\]
Подставляя и переписывая правую часть как интеграл от производной:
\[
  \int\limits_{x_1}^{x_2} u_t\, dx + a \int\limits_{x_1}^{x_2} u_x\, dx = 0 .
\]
Отрезок $[x_1,x_2]$ произволен, подынтегральное выражение непрерывно, поэтому
\[
  \boxed{\ u_t + a\, u_x = 0\ }
\]
— линейное однородное \textbf{уравнение переноса}. В обозначениях
задачи~\eqref{eq:fo-homog} направляющий вектор равен $\bar a = (a, 1)$ в
плоскости $(x,t)$.
\end{example}

\inputfig{figures/fig-transport-lines}

В $\R^n$ та же выкладка для среды, движущейся со скоростью $a$ в направлении
единичного вектора $\hat n$, даёт
\[
  u_t + a\, \dd{u}{\hat n} = 0
  \qquad\Longleftrightarrow\qquad
  u_t + \bar a \cdot \nabla u = 0, \qquad \bar a = a\,\hat n ,
\]
с направляющим вектором $(\bar a, 1)$ в пространстве-времени.

\subsubsection*{Геометрическая интерпретация}

Пусть $u$ — гладкое решение уравнения~\eqref{eq:fo-homog}. Тогда:

\begin{itemize}[leftmargin=1.6em]
  \item $u(x) = C$ — поверхность уровня;
  \item $\nabla u$ — нормаль к этой поверхности;
  \item условие $\bar a \cdot \nabla u = 0$ означает $\bar a \perp \nabla u$, то
        есть \textbf{поле направлений $\bar a$ касается поверхностей уровня}.
\end{itemize}

\begin{example}
Уравнение $u_x = 0$ для $u = u(x,y)$: здесь $\bar a = (1,0)$. Поверхности
(линии) уровня — прямые $y = \mathrm{const}$, и общее решение
\[
  u = F(y) ,
\]
то есть произвольная функция от $y$. Линия уровня $F(y) = C$.
\end{example}

\begin{idea}
Отсюда видна общая схема решения уравнения~\eqref{eq:fo-homog}: надо найти
функции, постоянные вдоль интегральных кривых поля $\bar a$ (их называют
\emph{первыми интегралами}), и общее решение будет произвольной функцией от
них. Поэтому следующий шаг — выписать ОДУ для этих кривых.
\end{idea}

\subsubsection*{Характеристическая система}

Интегральные кривые поля $\bar a$ задаются системой ОДУ

\stepcounter{equation}
\begin{equation}\tag{2}\label{eq:char-sys-prop}
  \frac{dx_1}{a_1(x)} = \dots = \frac{dx_n}{a_n(x)} .
\end{equation}
Вводя общее значение этих отношений как $ds$,
\[
  \frac{dx_1}{a_1(x)} = \dots = \frac{dx_n}{a_n(x)} = ds ,
\]
получаем

\stepcounter{equation}
\begin{equation}\tag{3}\label{eq:char-sys}
  \frac{dx_i}{ds} = a_i(x), \qquad i = \overline{1,n} ,
\end{equation}
с начальным условием

\stepcounter{equation}
\begin{equation}\tag{4}\label{eq:char-init}
  x_i(s_0) = x_i^{(0)}, \qquad i = \overline{1,n} ,
\end{equation}
или в векторной форме
\[
  \begin{cases}
    \dfrac{d\vec x}{ds} = \bar a(x), \\[8pt]
    \vec x(s_0) = \vec x_0 .
  \end{cases}
\]

\begin{idea}
Система~\eqref{eq:char-sys} \emph{автономна}: правая часть не зависит от $s$
явно. Это и даёт свойство, сформулированное ниже, — сдвиг параметра вдоль
решения снова даёт решение.
\end{idea}

\subsubsection*{Пример: гармонический осциллятор}

\begin{example}
Уравнение движения груза на пружине:
\[
  \begin{cases}
    m\ddot x = -\kappa x, \\
    x(t_0) = x_0, \quad \dot x(t_0) = v_0 .
  \end{cases}
\]
Сводим к системе первого порядка, вводя скорость $v$:
\[
  \begin{cases}
    \dot x = \dfrac{dx}{dt} = v, \\[6pt]
    m\dot v = m\dfrac{dv}{dt} = -\kappa x, \\[6pt]
    x(t_0) = x_0, \quad v(t_0) = v_0 .
  \end{cases}
\]
В обозначениях задачи~\eqref{eq:char-sys}--\eqref{eq:char-init} это значит
\[
  x_1 = x, \qquad x_2 = v, \qquad s = t ,
\]
а роль поля $\bar a$ играет
\[
  a_1(x,v) = v, \qquad a_2(x,v) = -\frac{\kappa x}{m} = -\omega^2 x,
  \qquad \omega^2 = \frac{\kappa}{m} .
\]
Решение:
\[
  x(t) = \alpha \sin(\omega t + \beta), \qquad
  v(t) = \alpha \omega \cos(\omega t + \beta) ,
\]
а траектории в фазовом пространстве — эллипсы
\[
  \omega^2 x^2 + v^2 = \alpha^2 \omega^2 .
\]
\end{example}

\subsubsection*{Траектории и интегральные кривые}

На этом примере удобно развести два понятия:

\begin{itemize}[leftmargin=1.6em]
  \item $(x(t), v(t)) \in \R^2$ — \textbf{фазовое пространство}, то есть
        пространство зависимых переменных;
  \item $(x(t), v(t), t) \in \R^3$ — \textbf{расширенное фазовое пространство};
  \item $(x(t), v(t), t)$ — \textbf{интегральная кривая} (график решения);
  \item $(x(t), v(t))$ — \textbf{траектория} (в фазовом пространстве).
\end{itemize}

\begin{proposition}
Пусть две интегральные кривые проходят через точки
\[
  (x_0, v_0, t_0) \qquad\text{и}\qquad (x_0, v_0, t_1) .
\]
Тогда они имеют одну и ту же траекторию.
\end{proposition}

В общем случае, для системы~\eqref{eq:char-sys}:

\begin{itemize}[leftmargin=1.6em]
  \item $\vec x(s)$ — траектория;
  \item $(\vec x(s), s) \in \R^{n+1}$ — расширенное фазовое пространство;
  \item $(\vec x(s), s)$ — интегральная кривая (график решения);
  \item интегральные кривые, проходящие через точки $(\vec x_0, s_0)$ и
        $(\vec x_0, s_1)$, имеют одну и ту же траекторию.
\end{itemize}

\begin{idea}
Причина — автономность: если $\vec x(s)$ решает~\eqref{eq:char-sys}, то
$\vec x(s + c)$ тоже решает её при любом $c$. Сдвиг по времени двигает
интегральную кривую вдоль оси $s$, но её проекция на фазовое пространство —
траектория — не меняется. Поэтому траектории не пересекаются и разбивают
фазовое пространство на непересекающиеся кривые, тогда как через одну точку
фазового пространства проходит целое однопараметрическое семейство
интегральных кривых.
\end{idea}

\begin{exercise}
Как выглядят интегральные кривые для гармонического осциллятора?
\end{exercise}

\inputfig{figures/fig-oscillator-phase}

\begin{idea}
Траектория — эллипс в плоскости $(x,v)$, проходимый с постоянной угловой
скоростью $\omega$. Добавляя ось $t$, получаем \textbf{эллиптическую винтовую
линию} в $\R^3$: проекция на $(x,v)$ — тот же эллипс, шаг по $t$ равен периоду
$T = 2\pi/\omega$. Разным начальным моментам $t_0$ отвечают винтовые линии,
сдвинутые друг относительно друга вдоль оси $t$, — ровно то, о чём утверждение
выше. Положению равновесия $x = v = 0$ отвечает вырожденная траектория-точка и
прямая, параллельная оси $t$.
\end{idea}

\subsubsection*{Траектории с общей точкой}

\begin{proposition}\label{prop:traj-coincide}
Траектории, имеющие общую точку, совпадают.
\end{proposition}

\begin{proof}
Пусть $\vec x_1(s)$ — решение~\eqref{eq:char-sys}--\eqref{eq:char-init} с
начальным условием $(\vec x_0, s_1)$, а $\vec x_2(s)$ — решение с начальным
условием $(\vec x_0, s_2)$. Поле $\bar a(x)$ гладкое, поэтому работает теорема
единственности для ОДУ. Утверждается, что
\[
  \vec x_2(s) = \vec x_1(s - s_2 + s_1) .
\]

Правая часть решает~\eqref{eq:char-sys} (система автономна, сдвиг
параметра снова даёт решение), а при $s = s_2$ она равна
$\vec x_1(s_1) = \vec x_0$. Значит она решает ту же задачу Коши, что и
$\vec x_2$, и по единственности обе функции совпадают. Кривые отличаются лишь
сдвигом параметра, то есть совпадают как траектории.
\end{proof}

\begin{idea}
Геометрический смысл: траектория — один и тот же объект, пройденный со сдвигом
параметра, и через каждую точку фазового пространства проходит ровно одна
траектория. Это то же самое утверждение, что и про проекции интегральных
кривых: кривые, сдвинутые вдоль оси $s$, проектируются в одну траекторию.
\end{idea}

\begin{definition}
Траектория системы~\eqref{eq:char-sys} называется её
\textbf{характеристикой}. Характеристиками уравнения~\eqref{eq:fo-homog}
называются характеристики системы~\eqref{eq:char-sys}.
\end{definition}

\begin{remark}
Характеристики системы~\eqref{eq:char-sys} описываются характеристической
системой~\eqref{eq:char-sys-prop}, в которой параметра $s$ уже нет: он нужен
был только чтобы перейти от пропорции к автономной системе, а сама кривая от
параметризации не зависит.
\end{remark}

\begin{warn}
Смазанная фраза на доске («хар-ки системы (3) описываются характеристическими
системами») — это именно замечание выше: в тетради она записана как «хар-ки
системы (3) описываются хар. системой (2) (нет переменной $s$)».
\end{warn}

\subsubsection*{Геометрическая интерпретация характеристик}

\begin{itemize}[leftmargin=1.6em]
  \item поле $\bar a$ в каждой точке касается характеристики;
  \item обратно, кривая, в каждой точке которой поле $\bar a$ касательно,
        является характеристикой.
\end{itemize}

Оба утверждения выводятся из пропорции~\eqref{eq:char-sys-prop} и теоремы
единственности для ОДУ: касание поля в каждой точке — это и есть
покоординатная пропорциональность $dx_i$ и $a_i(x)$, а дальше через каждую
точку проходит ровно одна такая кривая.

Пусть $u(\vec x)$ — гладкая функция. Тогда вдоль характеристики

\stepcounter{equation}
\begin{equation}\tag{5}\label{eq:char-deriv}
  \frac{d\,u(\vec x(s))}{ds}
  = \sum_{i=1}^{n} u_{x_i} \frac{dx_i}{ds}
  = \sum_{i=1}^{n} a_i(x)\, u_{x_i}
  = \bar a(x) \cdot \nabla u .
\end{equation}

\begin{remark}
Интерпретация соотношения~\eqref{eq:char-deriv}: \textbf{на характеристиках
оператор в частных производных $\bar a \cdot \nabla$ превращается в обыкновенную
производную по $s$}. Это и есть весь метод характеристик.
\end{remark}

\subsubsection*{Первые интегралы}

\begin{definition}
Функция $u(\vec x) \not\equiv \mathrm{const}$ называется \textbf{первым
интегралом} системы~\eqref{eq:char-sys}, если она постоянна на любом решении
этой системы.
\end{definition}

\begin{proposition}
Функция $u(x)$ является решением уравнения~\eqref{eq:fo-homog} тогда и только
тогда, когда она является первым интегралом системы~\eqref{eq:char-sys}.
\end{proposition}

\begin{proof}
Немедленно из~\eqref{eq:char-deriv}: равенство $\bar a \cdot \nabla u = 0$ в
точности означает, что $\frac{d}{ds} u(\vec x(s)) = 0$ вдоль каждой
характеристики.
\end{proof}

\begin{definition}
Первые интегралы $u_1,\dots,u_k$ называются \textbf{независимыми}, если
\[
  \operatorname{rank}
  \frac{\partial(u_1,\dots,u_k)}{\partial(x_1,\dots,x_n)} = k ,
\]
то есть не существует $F \not\equiv 0$, для которой
$F(u_1,\dots,u_k) \equiv 0$ в окрестности точки.
\end{definition}

\begin{warn}
Знаменатель якобиана: на доске он смазан, в тетради записан правильно —
$\partial(x_1,\dots,x_n)$, с пометкой «ранг максимален».

\medskip
А вот словесная расшифровка независимости и на доске, и в тетради записана
криво: «не существует функции $F(u_1,\dots,u_k) \not\equiv \mathrm{const}$ в
окрестности точки». По смыслу имеется в виду «не существует нетривиальной
гладкой $F$, для которой $F(u_1,\dots,u_k) \equiv \mathrm{const}$», то есть
функциональная независимость. В тетради рядом стоит более внятная приписка:
\emph{все градиенты не коллинеарны} — это ровно условие на ранг.
\end{warn}

\begin{theorem}\label{thm:first-integrals}
Система~\eqref{eq:char-sys} имеет ровно $n-1$ независимых первых интегралов;
обозначим их через $u_1,\dots,u_{n-1}$. Всякий первый интеграл имеет вид
\[
  F(u_1,\dots,u_{n-1}), \qquad F \not\equiv \mathrm{const} ,
\]
и обратно: для любой гладкой $F \not\equiv \mathrm{const}$ функция
$F(u_1,\dots,u_{n-1})$ является первым интегралом.
\end{theorem}

\begin{warn}
На доске обе половины утверждения остались недописанными; тетрадь их
подтверждает дословно: «всякий первый интеграл имеет вид
$F(u_1,\dots,u_{n-1})$, где $F$ гладкая и $F \not\equiv \mathrm{const}$, и
наоборот — всякая такая гладкая функция есть первый интеграл». Так что
восстановление верное.
\end{warn}

\begin{corollary*}
Общее решение уравнения~\eqref{eq:fo-homog} имеет вид
$u = F(u_1,\dots,u_{n-1})$.
\end{corollary*}

\begin{idea}
Почему именно $n-1$: траектории заполняют $\R^n$ однопараметрическими кривыми,
так что «поперёк» них остаётся $n-1$ независимых направлений. Первый интеграл —
это функция, постоянная вдоль траекторий, то есть функция на множестве самих
траекторий, а это множество $(n-1)$-мерно.
\end{idea}

\subsubsection*{Как искать первые интегралы}

Решаем систему~\eqref{eq:char-sys-prop} и получаем $n-1$ соотношений — это
решение в неявном виде:
\[
  \begin{cases}
    W_1(x,\, C_1,\dots,C_{n-1}) = 0, \\
    \ \ \vdots \\
    W_{n-1}(x,\, C_1,\dots,C_{n-1}) = 0,
  \end{cases}
\]
после чего разрешаем их относительно констант:
\[
  \begin{cases}
    u_1(x) = C_1, \\
    \ \ \vdots \\
    u_{n-1}(x) = C_{n-1} .
  \end{cases}
\]
Слева в этих равенствах и стоят искомые первые интегралы, а сама система при
произвольных $C_i$ задаёт \emph{произвольную} траекторию
системы~\eqref{eq:char-sys}.

\medskip
Чтобы получить \emph{конкретную} траекторию — ту, что проходит через точку
$x_0$, то есть решение задачи~\eqref{eq:char-sys}--\eqref{eq:char-init}, —
надо подставить вместо констант значения в этой точке:
\[
  \begin{cases}
    u_1(x) = u_1(x_0), \\
    \ \ \vdots \\
    u_{n-1}(x) = u_{n-1}(x_0) .
  \end{cases}
\]

\begin{idea}[title={Почему можно просто подставить $x_0$}]
Пометка «почему?» в тетради отсылает к системе~\eqref{eq:char-sys} с начальным
условием~\eqref{eq:char-init}. Вдоль траектории функции $u_i$ по определению
не меняются, а точка $x_0$ на этой траектории лежит. Значит константы обязаны
равняться $C_i = u_i(x_0)$ — другого выбора нет.
\end{idea}

\subsubsection*{Пример: уравнение переноса}

\begin{example}
Для $u = u(x,t)$ и уравнения $u_t + a u_x = 0$ характеристическая
система~\eqref{eq:char-sys-prop} выглядит как
\[
  \frac{dx}{a} = \frac{dt}{1} ,
\]
откуда
\[
  x - at = C \quad\text{— характеристики}, \qquad
  x - at \quad\text{— первый интеграл} .
\]
Общее решение:
\[
  u(x,t) = F(x - at) ,
\]
направляющий вектор $\bar a = (a,1)$ (см. рис.~\ref{fig:transport-lines}).

\medskip
Добавим начальное условие $u|_{t=0} = \varphi(x)$. Тогда $F(x) = \varphi(x)$, и
решение задачи равно
\[
  \boxed{\ u(x,t) = \varphi(x - at)\ }
\]
\end{example}

\begin{warn}
В тетради последняя строка не дописана: «$\Rightarrow$ реш-е задачи
$u(x,t) =$ \textbf{???}». Ответ восстановлен подстановкой $F = \varphi$.
\end{warn}

\begin{exercise}
Найти характеристики и общее решение уравнения
\[
  u_t + a u_x + b u_y = 0, \qquad u = u(x,y,t),
  \qquad \bar a = (a,\, b,\, 1) .
\]
\end{exercise}

\begin{idea}
Характеристическая система: $\dfrac{dx}{a} = \dfrac{dy}{b} = \dfrac{dt}{1}$.
Независимых первых интегралов здесь $n-1 = 2$ штуки: $x - at$ и $y - bt$.
Характеристики — прямые $x - at = C_1$, $y - bt = C_2$, а общее решение
$u = F(x - at,\ y - bt)$.
\end{idea}

\subsubsection*{Задача Коши на кривой}

Пусть данные заданы не на прямой $t=0$, а на кривой
\[
  \gamma = \big\{\, (\xi(\tau),\, \eta(\tau)) \,\big\},
  \qquad u|_\gamma = \varphi(\tau) .
\]

\begin{definition}
\textbf{Условие трансверсальности}: поле направлений $\bar a$ нигде не касается
$\gamma$.
\end{definition}

\inputfig{figures/fig-cauchy-gamma}

Подставляя общее решение $u = F(x - at)$ в начальное условие, получаем уравнение
на $F$:

\stepcounter{equation}
\begin{equation}\tag{6}\label{eq:transversal-eq}
  F\big(\xi(\tau) - a\,\eta(\tau)\big) = \varphi(\tau) .
\end{equation}

\begin{exercise}[теорема об обратной функции]
Докажите, что если $\xi, \eta, \varphi$ гладкие и выполнено условие
трансверсальности, то уравнение~\eqref{eq:transversal-eq} можно однозначно
разрешить относительно $F$, и значит решение задачи Коши существует и
единственно в окрестности $\gamma$.

\medskip
\emph{Подсказка:} рассмотреть $g(\tau) = \xi(\tau) - a\,\eta(\tau)$.
\end{exercise}

\begin{idea}
Разбор подсказки. Касательный вектор к $\gamma$ равен $(\xi', \eta')$, а
направляющий вектор характеристики — $\bar a = (a, 1)$. Они параллельны тогда и
только тогда, когда
\[
  \det \begin{pmatrix} \xi' & a \\ \eta' & 1 \end{pmatrix}
  = \xi'(\tau) - a\,\eta'(\tau) = 0 .
\]
Но это в точности $g'(\tau)$. Значит условие трансверсальности равносильно
$g'(\tau) \ne 0$: функция $g$ строго монотонна и потому локально обратима.
Тогда~\eqref{eq:transversal-eq} даёт
\[
  F = \varphi \circ g^{-1} ,
\]
то есть $F$ восстанавливается однозначно, а решение равно
$u(x,t) = \varphi\big(g^{-1}(x - at)\big)$. Гладкость $F$ наследуется от
$\varphi$ и $g^{-1}$.

\medskip
Короче говоря: \textbf{трансверсальность $\iff$ якобиан отличен от нуля}, и всё
упражнение — это теорема об обратной функции.
\end{idea}

\subsubsection*{Что бывает без трансверсальности}

\begin{example}
\[
  \begin{cases}
    u_x = 0, \\
    u|_{y = x^3} = x^2 .
  \end{cases}
\]
Здесь $\bar a = (1,0)$, характеристики — прямые $y = \mathrm{const}$, а кривая
$y = x^3$ касается одной из них в начале координат: её касательная
$(1, 3x^2)$ при $x = 0$ равна $(1,0)$.

\medskip
Общее решение $u = F(y)$; из начального условия $F(x^3) = x^2$, то есть
\[
  u = y^{2/3} ,
\]
и это функция, \textbf{негладкая на оси} $Ox$.
\end{example}

\inputfig{figures/fig-tangency}

\begin{idea}
Трансверсальность нарушена ровно в одной точке — и этого уже достаточно, чтобы
решение потеряло гладкость. Причина видна по формуле $F = \varphi \circ g^{-1}$
из предыдущего пункта: в точке касания $g' = 0$, обратная функция $g^{-1}$
негладкая, и эта негладкость переносится на $F$. Здесь $g(x) = x^3$,
$g^{-1}(y) = y^{1/3}$, $\varphi(x) = x^2$ — и получается ровно $y^{2/3}$.
\end{idea}
